% Imporditud: tools/import_eksperimendid.py
% Allikas: Eesti füüsikaolümpiaadi eksperimendi ülesannete
%   kogu (Rael Kalda, 2023), ülesanne Ü140, lahendus L140.
\setAuthor{EFO žürii}
\setRound{piirkonnavoor}
\setYear{2011}
\setNumber{G E1}
\setDifficulty{4}
\setTopic{geomeetriline optika}
\setVahendid{luup, mõõtejoonlaud, millimeetripaber}
\setPunktid{10}
\setAllikas{Eksperimendi kogu Ü140, lahendus lk 105}
\setLahendusseis{toores}

\prob{Luubi suurendus}
Tehke kindlaks kuidas oleneb luubi suurendus kaugusest luubi ja eseme vahel, kui kaugus eseme ja silma vahel hoida $\approx$ 30 cm. Tulemus esitage graafiliselt.
\textit{Märkus:} Suurendus näitab, mitu korda suureneb luubi abil vaadeldava väikese eseme nurkläbimõõt.
\textit{Vihje:} Väikese eseme fikseeritud kauguse korral on selle nurkläbimõõt võrdeline joonmõõtmega.

\hint

\solu
Asetame luubi millimeetripaberist erinevatele kaugustele ja võrdleme erinevatel kaugustel ruudustiku suurust läbi luubi ja ilma luubita. Mõõdame luubi kaugused millimeetripaberist ja hindame vastavad suurendused. Joonestame graafiku, millel on antud suurenduse sõltuvus eseme kaugusest luubist, kanname graafiku väljale punktid ja ühendame need. Teoreetiliselt tuletatud avaldis luubi suurenduse jaoks (vastavalt tekstis toodud definitsioonile):

fl S = l(f $-$ a) + $a_2$ ,

kus l on silma ja objekti vahekaugus ning a --- objekti ja läätse vahekaugus. Kui f > l/2, siis omab suurendus kui funktsioon a-st maksimumi a = l/2 juures,

Sm = f . l/4 f $-$

Väärtuste l = 30 cm ja f = 18 cm jaoks on graafik toodud juuresoleval joonisel.

Hindamine: Idee katse teostamiseks --- 3 p. 4-5 erineval kaugusel mõõtmist --- 2 p. Mõõtmiste põhjal suurenduste arvutamine --- 2 p. Tulemuste arvväärtused on mõistlikud ($\pm$30\%) --- 1 p. Korrektne graafik --- 2 p.
\probend
